% SPDX-FileCopyrightText: 2026 Yiyu Zhou <yiyu@yiyuzhou.io>
%
% SPDX-License-Identifier: GPL-3.0-or-later

\chapter{Keys}
\label{ch:keys}

\section{Keymaps}

cinq ships one keymap, vi.  It moves with letters such as \key{j} and
\key{k}, and uses two-key sequences for some actions, such as \key{dd}.  It
opens the help page with \key{?}, which departs from Vim, and with \key{F1}.
It puts GO on \key{Space} and Panic on \key{Escape}, as QLab does, and also
uses \key{ArrowUp}, \key{ArrowDown}, \key{Enter}, \key{Alt} with those
arrows, \key{Control+s}, and \key{F11}.

The actions that this tutorial performed with buttons and clicks also have
keys:

\begin{center}
\small
\begin{tabular}{@{}p{0.62\textwidth}p{0.34\textwidth}@{}}
\hline
Action & vi \\
\hline
Add a cue below the selection & \key{o} \\
Cut the selected cue & \key{dd}, \key{x}, \key{Delete} \\
Copy the selected cue & \key{yy}, \key{Y} \\
Paste below the selection & \key{p} \\
Paste above the selection & \key{P} \\
Move the selected cue into the group above it & \key{>>} \\
Move the selected cue out of its group & \key{<<} \\
Cancel what is typed halfway & \key{Control+c}, \key{Control+[} \\
Edit the name of the selected cue & \key{i}, \key{a}, \key{A} \\
Edit the name from its start & \key{I} \\
Change the name, selecting all of it & \key{cc}, \key{S} \\
Undo the latest change & \key{u} \\
Redo the change undone last & \key{Control+r} \\
Open a show & \key{Control+l} \\
Put the selected cue on standby & \key{Enter} \\
Show the next cue list, then the outputs & \key{l} \\
Show the previous cue list & \key{h} \\
Move down half a page & \key{Control+d} \\
Move down a page & \key{Control+f} \\
Show the next inspector tab & \key{gt} \\
Show inspector tab number 2, or another & \key{2gt} \\
Go back to the place before the last jump, or leave the outputs & \key{Control+o} \\
Go forward again & \key{Control+i} \\
Search the cues, or the outputs, by their text & \key{/} \\
Select the next match & \key{n} \\
Show or hide the help page & \key{?} or \key{F1} \\
Pick the highlighted option of an open menu & \key{Control+y} \\
Close an open menu without picking & \key{Control+e} \\
\hline
\end{tabular}
\end{center}

\key{Enter} picks as well.  While a drop-down list is open, the status
bar names the keys that pick and close it, and \key{PageDown}, \key{PageUp},
\key{Home}, and \key{End} move the highlight by a page or to either end.

\section{The help page}

\begin{enumerate}
\item Press the help key that the status bar names, \key{?} or \key{F1}.
  The help page opens with a paragraph of first steps, then lists
  every action that has a key, with its keys and, in a wide window, its name
  in the configuration file, as Figure~\ref{fig:help} shows.
\item Press \key{Enter} on the page to open the tutorial, which the first
  paragraph names.
\item Press the same key again to close the page, or \key{q}, which only
  closes it.  GO and Panic also close it, so that the show stays in view
  while it runs.
\end{enumerate}

\begin{figure}[H]
  \centering
  \includegraphics[width=\textwidth]{screenshots/help.png}
  \caption{Help page of the vi keymap}
  \label{fig:help}
\end{figure}

\section{Key sequences}

Some actions take more than one key, such as \key{gg}, which selects the first
cue.  The keys of a sequence are typed one after another with no \key{Space}
between them, since \key{Space} is GO.

\begin{enumerate}
\item Press \key{g}.  The status bar shows the keys typed so far, and a box
  lists the keys that can follow, as in Figure~\ref{fig:keys}.  While the
  outputs show, it names outputs where it names cues here.
\item Press \key{g} again.  The first cue is selected.
\end{enumerate}

\begin{figure}[H]
  \centering
  \includegraphics[width=\textwidth]{screenshots/keys.png}
  \caption{Keys that can follow \key{g}}
  \label{fig:keys}
\end{figure}

A number typed first repeats the action: \key{3j} moves the selection three
rows down, \key{3dd} cuts three cues, \key{3o} adds three cues, \key{3p}
pastes the copied cues three times, and \key{3u} undoes three changes.
\key{d} and \key{y} also take a motion: \key{dj} cuts the selected cue and
the next, \key{yk} copies it and the previous, and \key{dG} cuts down to the
last cue.  A number before the operator, before the motion, or both
multiplies, so that \key{2d3j} cuts seven cues.  A key that does not
continue the sequence cancels it and does nothing: cinq drops the keys typed
so far without a sound or a message, so that a mistyped \key{dx} cuts nothing
and the status bar goes back to its usual hints.  The keys for GO and Panic
act immediately, even in the middle of a sequence.  \key{Control+c} cancels a
sequence, as it does in Vim, and the status bar hints it while a sequence
waits or a text field has the keyboard, so that no one reaches for
\key{Escape}, which is Panic.

\section{Searching the cues}

A long cue list is easier to search than to scroll.

\begin{enumerate}
\item Press the search key, \key{/}.  A search line opens at the left of
  the status bar.
\item Type part of what the list shows of a cue, such as its name, its number,
  or the file that it plays.  The selection moves to the first cue after the
  selection that holds the text, in any case, and the list scrolls to it.  The
  search wraps around the list and ends on the selection, so that the
  selection is the match only when no other cue holds the text.  When no cue
  holds the text, the line says so and the selection goes back to where the
  search began.
\item While typing, \key{Control+w} deletes the word before the caret and
  \key{Control+u} everything before it, as in any text field, and
  \key{ArrowUp} and \key{ArrowDown} recall earlier searches that start with
  the text.
\item Press \key{Enter} to accept the search.  \key{Control+c} and
  \key{Escape}, which is also Panic, close the line and put the selection
  back, and any click in the main window closes it and leaves the selection
  where it is.
\item Press \key{n} for the next match in the direction of the search, wrapping
  around at the end of the list, and \key{N} for the next match in the other
  direction.  \key{Enter} on an empty line repeats the last search.
\end{enumerate}

Searches are jumps, as in Vim.  After a search, \key{Control+o} goes back to
where it began, or to the cue before a step to the next match, and
\key{Control+i} returns to the match.

\section{Hint mode}

The keys do not reach everything that a click does, so a key can label each
place for you, as Vimium does in a browser.

\begin{enumerate}
\item Press the hint key, \key{f}.  A small label of a few
  letters appears left of each button, check box, inspector tab, cue list of
  the sidebar, drop-down list, text field, and group mark that is in view.
  The status bar asks for a label.
\item Type the letters of one label.  The labels that the letters do not begin
  disappear, and \key{Backspace} takes a letter back.  When the letters are
  whole, cinq acts as a click on that place would: it presses the button,
  flips the check box, shows the tab or the cue list, or gives the text field
  or the drop-down list the keyboard, where \key{Enter} opens the menu.
\item To leave without choosing, press \key{Control+c}, or \key{Backspace}
  where no letter is typed, or click in the main window.  A letter that begins
  no label is ignored, so that a slip never runs the next letters as commands.
  GO and Panic act as ever and leave too, and so does a cue that starts or ends
  in the Active list of the sidebar, since the labels would no longer lie on
  their places.
\end{enumerate}

The help page takes labels on its links.  The rows of the cue list have their
own keys.  GO and Panic have no label, as their keys act from everywhere.

\section{Changing the keys}

The configuration file, \texttt{config.toml}, binds actions to other keys and
declares keymaps of its own.  The help page links to it, and following the link
creates the file if it does not exist.  cinq rereads the file whenever it
changes, so edits take effect as soon as the file is saved.  The manual page
\texttt{cinq(5)} describes every setting, lists the keys of the vi keymap, and
shows how to declare a keymap of one's own.

\section*{Exercise}

You can now find the key for any action and choose the keys that suit you.
Declare a keymap of your own that starts from vi, as the examples in
\texttt{cinq(5)} show, and open the help page with \key{?}.
